\documentclass[aspectratio=169, 14pt]{beamer}
\usepackage{beamer_preamble}

\begin{document}

% ==== Title ===============================================================
\titleslide{Lesson 4-1 \textperiodcentered\ Period 2}{Inverse Variation Applications + DOK-3 Performance Task}

% ==== Do Now ==============================================================
\begin{frame}{Do Now --- Is this inverse variation?}{Algebra 2 \textperiodcentered\ Lesson 4-1 \textperiodcentered\ Period 2}
\phasebadges{1}{5}{bridge from P1}
\vspace{0.2cm}
\begin{projcard}{Practice \#15 \textperiodcentered\ Table check}{slidegold}
\begin{tabular}{l|cccccc}
\(x\) & 1 & 2 & 3 & 4 & 5 & 6 \\
\(y\) & 60 & 30 & 20 & 15 & 12 & 10 \\
\end{tabular}

\vspace{0.25cm}
Is this an inverse variation? Compute \(xy\) for each column.
If yes, state \(k\).
\end{projcard}
\end{frame}

% ==== Launch ==============================================================
\begin{frame}{Launch --- Example 3 (Bouzouki)}{Algebra 2 \textperiodcentered\ Lesson 4-1 \textperiodcentered\ Period 2}
\phasebadges{2}{12}{teacher models}
\vspace{0.2cm}
\begin{projcard}{Application Rules for Inverse Variation}{slidegreen}
\begin{enumerate}
  \setlength{\itemsep}{0pt}
  \item Identify two varying quantities and their constant PRODUCT \(k\).
  \item Compute \(k\) from one given pair.
  \item Write the equation \(y = k/x\) (rename variables to fit).
  \item Solve for the unknown.
\end{enumerate}

\vspace{0.1cm}
\small
Example (bouzouki string): length \(s\) varies inversely with frequency \(f\).\\
\hspace*{1em}\(\bullet\) 26-inch string \(\rightarrow\) 329.63 cycles/sec \(\rightarrow\) \(k = 8{,}570.38\)\\
\hspace*{1em}\(\bullet\) 13-inch string \(\rightarrow\) \(f = 8{,}570.38 / 13 = 659.26\) cycles/sec

\vspace{0.1cm}
\textcolor{slideaccent}{\textbf{DIRECT (\(y/x = k\)) vs INVERSE (\(xy = k\))} --- both appear in today's DOK-3!}
\end{projcard}
\end{frame}

% ==== Explore =============================================================
\begin{frame}{Explore --- TPS + DOK-3 Performance Task}{Algebra 2 \textperiodcentered\ Lesson 4-1 \textperiodcentered\ Period 2}
\phasebadges{2-3}{33}{student work}
\small\textcolor{slidegray}{5 items in order. Plan \(\sim\)15 min for Practice \#26 (Ramón).}

\vspace{0.15cm}
\begin{projcard}{Try It 3 \textperiodcentered\ Ice cube melting}{slideblue}
Time to melt varies inversely with air temperature.
\end{projcard}
\vspace{-0.15cm}
\begin{projcard}{Practice \#12 \textperiodcentered\ HOT: inverse SQUARE}{slideblue}
\(y = k/x^2\). If \(x\) is multiplied by 4, what happens to \(y\)?
\end{projcard}
\vspace{-0.15cm}
\begin{projcard}{Practice \#17 \textperiodcentered\ Radio wavelength (graph)}{slideblue}
\(w\) varies inversely with frequency \(f\). Find \(f\) when \(w = 375\) m.
\end{projcard}
\vspace{-0.15cm}
\begin{projcard}{Practice \#23 \textperiodcentered\ Boyle's Law}{slideblue}
Volleyball 300 in\(^3\) at 4.5 psi vs basketball 415 in\(^3\) at 8 psi.
\end{projcard}
\vspace{-0.15cm}
\begin{projcard}{Practice \#26 \textperiodcentered\ PERFORMANCE TASK (\(\bigstar\))}{slideaccent}
Ramón's road trip --- Part A is \textbf{DIRECT}, Part B is \textbf{INVERSE}.
\end{projcard}
\end{frame}

% ==== Share / Summary =====================================================
\begin{frame}{Share / Summary}{Algebra 2 \textperiodcentered\ Lesson 4-1 \textperiodcentered\ Period 2}
\phasebadges{2}{5}{DOK-3 reveal + Wed bridge}
\vspace{0.2cm}
\begin{projcard}{Sentence frame \textperiodcentered\ Ramón}{slideblue}
``Part A is \underline{\hspace{1cm}} variation because distance and time
have a constant \textbf{RATIO} of \underline{\hspace{1cm}}.\\
Part B is \underline{\hspace{1cm}} variation because gas and time have
a constant \textbf{PRODUCT} of \underline{\hspace{1cm}}.''
\end{projcard}

\vspace{0.15cm}
\begin{projcard}{Bridge to Period 3 (Wednesday)}{slideaccent}
Tomorrow we graph \(y = 1/x\) and its translations.
Watch for: asymptotes, domain, range, intercepts.
Period 3 is \textbf{assessment-critical} for the Topic 4 8-Q assessment.
\end{projcard}
\end{frame}

% ==== Exit Ticket =========================================================
\begin{frame}{Exit Ticket --- Summary Recap}{Algebra 2 \textperiodcentered\ Lesson 4-1 \textperiodcentered\ Period 2}
\phasebadges{1-2}{3}{not graded}
\vspace{0.2cm}
\begin{projcard}{Complete on your exit ticket}{slidegold}
\begin{enumerate}
  \setlength{\itemsep}{0.2cm}
  \item Today I learned that \underline{\hspace{5cm}}.
  \item The biggest trap in Ramón's task was \underline{\hspace{5cm}}.
  \item I distinguish direct from inverse variation by \underline{\hspace{5cm}}.
\end{enumerate}
\end{projcard}
\end{frame}

\end{document}
